\section{Complementary faces control sharing between source rows}
\label{sec:face-sharing}
Unlike Section~\ref{sec:whole-rows}, we now allow every row to split
over several blocks, but require globally $C^1$ factor maps for the full
slack. Then the exposed-face dimensions of a block limit how many source
rows can be active in it, that is, have a nonzero mixed derivative there.
Globally, the exact Lorentz bounds of Theorem~\ref{thm:smooth-lorentz}
extend to every cone whose nonzero proper exposed faces are rays, even
with nonsmooth boundaries.

Let $M=(\Sph^p)^k$, $p=s-1\geq2$, and suppose
\begin{equation}\label{eq:proper-full-rows}
 1-\ip{x_a}{z}=\sum_{i=1}^L\ip{A_i(x)}{B_i^a(z)},
 \qquad x\in M,\ z\in\Sph^p,\ a\in[k],
\end{equation}
where $A_i:M\to K_i$ and $B_i^a:\Sph^p\to K_i^*$ are $C^1$ and the
cones are proper of dimensions $3\leq m_i\leq d$, for an integer
$d\geq3$. Set
\[
 r_i=m_i-2,\quad c=d-2,\quad
 f_i=\max_{\substack{0\ne b\in\partial K_i^*\\K_i\cap b^\perp\ne\{0\}}}
       \dim(K_i\cap b^\perp),\qquad \bar f_i=\min(f_i,k).
\]
We call $r_i$ the capacity of block $i$; block $i$ is active on row $a$
at $x$ when $t_i^a(x)>0$, with $t_i^a$ as in
Theorem~\ref{thm:face-sharing-local}. Here $f_i$ is a maximum of face
dimensions, not of numbers of extreme rays; below it may be replaced by
any upper bound on the dimensions of the nonzero exposed faces met by
the dual factors. Ray factors may be added: their mixed forms
vanish, so they contribute nothing to the bounds.

\begin{theorem}[Local face packing]\label{thm:face-sharing-local}
At $x\in M$ put
\[
 H_i^a(u,v)=-\ip{D_aA_i(x)[u]}{DB_i^a(x_a)[v]},\qquad
 t_i^a=\rank H_i^a,\qquad I_i=\{a:t_i^a>0\}.
\]
Then $t_i^a\leq r_i$, $\sum_i t_i^a\geq p$, and
\begin{equation}\label{eq:face-packing}
 \sum_{a\in I_i}t_i^a\leq r_i,\qquad
 1+\sum_{b\in I_i\setminus\{a\}}t_i^b
 \leq\dim(K_i\cap B_i^a(x_a)^\perp)\quad(a\in I_i).
\end{equation}
In particular $|I_i|\leq\bar f_i$.
\end{theorem}
\begin{proof}
Contact and nonnegativity make every pairing
$\ip{A_i(x)}{B_i^a(x_a)}$ zero. On the cylinder $x_a=u$, the function
$z\mapsto\ip{A_i(x)}{B_i^a(z)}$ has a minimum at $u$. Consequently
$\ip{A_i(x)}{DB_i^a(u)[v]}=0$ on that entire cylinder. Differentiating
in a different source direction, with the dual derivative held fixed,
gives
\[
 \ip{D_bA_i(x)[w]}{DB_i^a(x_a)[v]}=0\quad(b\ne a).
\]
This uses only first derivatives. The ray-quotient contact argument of
Lemma~\ref{lem:universal-channel} gives $t_i^a\leq m_i-2$; at a zero factor
its derivative is zero by pointedness. Mixed differentiation of the
full row gives $\sum_iH_i^a=g_{\Sph^p}$, hence $\sum_i t_i^a\geq p$.

Fix one block with an active row, write $A=A_i(x)\ne0$, and let
$\Lambda_a(v)$ be the functional $w\mapsto\ip{v}{DB_i^a(x_a)[w]}$.
Choose $E_a\subset\operatorname{im}D_aA_i$ of dimension $t_i^a$ on
which $\Lambda_a$ is injective. The identities just proved imply
$\Lambda_aA=0$ and $\Lambda_aE_b=0$ for $b\ne a$. Applying each
$\Lambda_a$ to a linear dependence proves that
\[
 \R A\ \oplus\!\bigoplus_{a\in I_i}E_a
\]
is a direct sum. Fix any active row $a_0$ and put
$b=B_i^{a_0}(x_{a_0})\ne0$. The nonnegative function
$x'\mapsto\ip{A_i(x')}{b}$ vanishes at $x$, so
$\operatorname{im}DA_i(x)\subset b^\perp$. Since $A\in b^\perp$,
the entire direct sum lies in this hyperplane of dimension $m_i-1$.
Thus $1+\sum_a t_i^a\leq m_i-1$, proving the first inequality.
When there is no active row that inequality is immediate. On the contact
cylinder for row $a$, the primal factor lies in
$F_a=K_i\cap B_i^a(x_a)^\perp$. Thus $E_b\subset\spanop F_a$ for
$b\ne a$, and $A\in F_a$. Their direct sum proves the second
inequality. Since each active rank is at least one, it also gives
$|I_i|\leq f_i$.
\end{proof}

\begin{theorem}[Global face budgets]\label{thm:face-sharing-global}
Define
\[
 h=\begin{cases}\lceil(p+1)/c\rceil,&c<p,\\1,&c\geq p,\end{cases}
 \quad
 \tau=\begin{cases}p+1,&c<p,\\p,&c\geq p,\end{cases}
 \quad R=\sum_i r_i,\qquad K=kp+\mathbf1_{\{c<p\}}.
\]
Every factorization~\eqref{eq:proper-full-rows} satisfies
\begin{equation}\label{eq:face-global-budgets}
 \sum_i\bar f_i\geq kh,\qquad
 \sum_i\bar f_i r_i\geq k\tau,\qquad R\geq K.
\end{equation}
If $f_i\leq f$ for a common $f\geq1$, set
\[
 R_*=\max\{\lceil k\tau/f\rceil,K\},\qquad
 L_*=\max\{\lceil kh/f\rceil,\lceil R_*/c\rceil\}.
\]
Then
\begin{equation}\label{eq:face-global-dimension}
 L\geq L_*,\qquad R\geq R_*,\qquad
 D:=\sum_i m_i\geq R_*+2L_*.
\end{equation}
Every self-concordant barrier on the ambient product has parameter
at least $2L_*$, even if it couples the factors or is not
logarithmically homogeneous.
\end{theorem}
\begin{proof}
Write $N_a(x)=|\{i:t_i^a(x)>0\}|$ and
$C_a(x)=\sum_{i:t_i^a(x)>0}r_i$. Pointwise,
$N_a\geq\lceil p/c\rceil$ and $C_a\geq p$. Counting incidences gives
\[
 \sum_aN_a\leq\sum_i\bar f_i,\qquad
 \sum_aC_a\leq\sum_i\bar f_i r_i.
\]
Summing~\eqref{eq:face-packing} gives $kp\leq R$. To exclude equality
when $c<p$, aggregate the labelled dual maps:
\[
 \widehat B_i(z_1,\ldots,z_k)=\sum_a B_i^a(z_a),\qquad
 \sum_i\ip{A_i(x)}{\widehat B_i(z)}
       =k-\sum_a\ip{x_a}{z_a}.
\]
These are $C^1$ cone-valued maps on $M$, with contact map $z=x$ and
nondegenerate mixed pairing equal to the product spherical metric.
At $R=kp$, Theorem~\ref{thm:general-joint-cover} therefore forces the
positive capacities to be exactly $k$ copies of $p$. This contradicts
$r_i\leq c<p$. Integrality gives $R\geq kp+1$ in this case.
Zero-capacity ray factors do not affect this argument. These observations
prove the last bound in~\eqref{eq:face-global-budgets}, and all its
bounds when $c\geq p$.

Suppose $c<p$. If $\sum_i\bar f_i r_i<k(p+1)$, every point belongs to
some $Z_a=\{C_a=p\}$. Each $Z_a$ is closed because activity is open
and $C_a\geq p$ everywhere. For a label set $J$, let $E_{a,J}$ be the
set of points of $Z_a$ where the active labels for row $a$ are exactly
$J$. These finitely many sets are clopen in $Z_a$, hence compact: an
active label stays active nearby, while a new one would raise $C_a$
above $p$. On such a piece
$\sum_{i\in J}r_i=p$. Normalize the nonzero dual factors $B_i^a$
using an interior primal vector to obtain maps into the boundaries of
compact dual-cone bases. These maps are defined near the projection
of $E_{a,J}$ onto its source sphere. Indeed fix the other source
coordinates at any point of the piece; nearby contact primal factors
remain nonzero and force the corresponding dual factors to remain on
the nonzero boundary.

The joint derivative of these normalized dual maps is injective at
every projected point. If a tangent vector is killed by all normalized
derivatives, each unnormalized derivative is radial and therefore
annihilates every primal contact derivative. The identity
$\sum_iH_i^a=g_{\Sph^p}$ then kills that tangent vector itself.
Injectivity persists on an open neighborhood. Lemma~\ref{lem:proper-base-domain}
shows that this neighborhood is a proper open subset of $\Sph^p$,
since every $r_i<p$ and their sum is $p$.

For a fixed $a$, choose pairwise disjoint open neighborhoods of the
finitely many compact pieces $E_{a,J}$, each inside the inverse image
of its proper source neighborhood. The pulled-back top class $u_a$ of
the $a$th source sphere vanishes on their union
$U_a$. The $U_a$ cover $M$, so the relative
cup-product argument of Section~\ref{sec:regularity} contradicts
$u_1\cdots u_k\ne0$. Thus $\sum_i\bar f_ir_i\geq k(p+1)$.
For the bound on $\sum_i\bar f_i$, if $c\nmid p$ the pointwise count already equals
$h$. Otherwise write $p=tc$. If $\sum_i\bar f_i<k(t+1)$, then $M$
is covered by the loci where some source has exactly $t$ active
labels. Every label there must have capacity $c$, and the same
compact-piece argument applies. This proves the count bound.

Uniform face caps and $R\geq K$ imply $R\geq R_*$ and
$L\geq\lceil kh/f\rceil$. Since $R\leq cL$, they also imply
$L\geq\lceil R_*/c\rceil$. The dimension bound follows from
$D=R+2L$. Finally each proper cone of dimension at least two
has a two-dimensional linear section through an interior point that
is a pointed wedge. Its relative boundary lies in the cone boundary.
The product sections form a $2L$-orthant. Restriction of any ambient
barrier and the orthant lower bound give $\nu\geq2L\geq2L_*$.
\end{proof}

\begin{corollary}[Exact minima for ray-exposed cones]\label{cor:ray-exposed-frontier}
Suppose every nonzero proper exposed face of each $K_i$ is a ray.
Over all dictionaries of such cones of dimension at most $d$, the minima
are attained simultaneously and equal
\[
 L=kh,\qquad R=k\tau,\qquad D=k\tau+2kh,
 \qquad \nu_{\rm ambient}=2kh,
\]
where $\nu_{\rm ambient}$ is the least barrier parameter on the
ambient product. They are attained by Lorentz cones and the grouped constructions of
Theorem~\ref{thm:smooth-lorentz}. A fixed dictionary need not contain
these attaining cones. We do not claim exact values for $f>1$.
\end{corollary}

The face inequalities also limit a single shared block. If block $i$ is
active on $q_i\geq2$
rows and $T_i=\sum_at_i^a$, summing its face inequalities gives
\begin{equation}\label{eq:face-sharing-tradeoff}
 (q_i-1)T_i\leq q_i(f_i-1),\qquad
 T_i\leq\min\{r_i,2(f_i-1)\}.
\end{equation}
Thus a small exposed face can further restrict a block's total mixed
rank when it serves several rows; sharing never increases that rank
beyond its capacity $r_i$. For a real PSD factor whose contact dual
has nullity $h_{ia}$, the face inequality reads
$1+\sum_{b\ne a}t_i^b\leq h_{ia}(h_{ia}+1)/2$; the sharper one-row
bound of Lemma~\ref{lem:peirce-channel} holds at the same time.

\subsection{A sharp model for shared scalar channels}
The next cone shows how split rows can share one block. Define
\[
 P_q=\{(t,u,w)\in\R\times\R^q\times\R^q:
             t\geq0,\ w_a\geq0,\ tw_a\geq u_a^2\}.
\]
This closed proper cone has dimension $2q+1$. The polynomial factors
\begin{equation}\label{eq:shared-square-model}
 A(v)=(1,v,v^{\circ2}),\qquad
 B^a(z)=(z^2/2,-ze_a,e_a/2)
\end{equation}
give $\ip{A(v)}{B^a(z)}=(v_a-z)^2/2$ on $[-1,1]^q\times[-1,1]$.
Dual feasibility follows directly by completing a square when $t>0$
and by continuity when $t=0$.

\begin{proposition}[The shared-square cone]\label{prop:shared-square}
The cone $P_q$ has maximum nonzero proper exposed-face dimension
$2q-1$. Its dimension $2q+1$ and this face dimension are the smallest possible
for a one-block factorization of all the kernels
in~\eqref{eq:shared-square-model}.
Its optimal self-concordant barrier parameter is $q+1$, attained by
\begin{equation}\label{eq:shared-square-barrier}
 F_q(t,u,w)=-\sum_{a=1}^q\log(tw_a-u_a^2)+(q-1)\log t.
\end{equation}
For a partition of $k$ source balls into $g$ groups of sizes $q_\ell$,
using these cones coordinate by coordinate gives an analytic
factorization of the full slack with $sg$ factors of total dimension $s(2k+g)$, maximum block dimension
$2\max q_\ell+1$, maximum exposed-face dimension $2\max q_\ell-1$,
and exact ambient parameter $s(k+g)$.
\end{proposition}
\begin{proof}
The dual cone consists of $\gamma\geq0$, $\beta_a=0$ when
$\gamma_a=0$, and
$\alpha\geq\sum_{\gamma_a>0}\beta_a^2/(4\gamma_a)$.
Let $J=\{a:\gamma_a>0\}$. In the equality case for a nonzero dual
vector, $J$ is nonempty and the contact equations fix
$u_a=-t\beta_a/(2\gamma_a)$ and
$w_a=t\beta_a^2/(4\gamma_a^2)$ for $a\in J$.
The other coordinates range over a copy of $P_{q-|J|}$, so the face
dimension is $1+2(q-|J|)$. In the strict case $t=0$, hence $u=0$;
only $w_a$ for $a\notin J$ remain, giving dimension $q-|J|$.
These cases give the maximum $2q-1$, including $q=1$.

The row functions span the linearly independent family
$1,v_a,v_a^2$ of size $2q+1$, proving the dimension lower bound.
On a contact cylinder $v_a=z$, the other rows linearly recover
$1,v_b,v_b^2$ for $b\ne a$; a nonzero constant is also recovered
from a different dual argument in row $a$. The contact face must
therefore span at least $1+2(q-1)$ dimensions.

The cone $P_q$ is the PSD-completable partial-matrix cone on a star:
the center diagonal is $t$, its edge entries are $u_a$, and the leaf
diagonals are $w_a$. For $t>0$ its maximum-determinant completion
has determinant $t\prod_a(w_a-u_a^2/t)$. The classical completion
barrier proved in Section~\ref{sec:spectral-chordal} consequently
is~\eqref{eq:shared-square-barrier}, of parameter $q+1$.
Equivalently it is the Legendre barrier of the sparse arrow-matrix
log determinant. Restricting any barrier to $u=0$ gives a
$(q+1)$-orthant with boundary on the original boundary, so the
parameter is optimal against arbitrary barriers.

For source $a$ in a group use $v_b=x_{b,j}$ and $z=z_j$ in coordinate
$j$. Summing gives $\tfrac12\sum_j(x_{a,j}-z_j)^2=1-\ip{x_a}{z}$
on the unit spheres. Summing dimensions and parameters gives the
stated counts; the product orthant sections prove the ambient lower
bound even for coupled barriers.
\end{proof}

For comparison, the rank-one real PSD factors
$A(v)=(1,v)(1,v)^T$ and
$B^a(z)=(e_a-ze_0)(e_a-ze_0)^T/2$ give the same scalar kernels,
but expose a PSD face of dimension $q(q+1)/2$. A whole-row version
uses $v=(1,x_1,\ldots,x_q)$ and sums those dual outer products over
the $s$ coordinates of a row; its dual rank is $s$, its nullity is
$1+(q-1)s$, and its contact face has dimension
$\binom{2+(q-1)s}{2}$. These PSD constructions are not claimed to be
optimal among homogeneous cones.

The homogenization of a group of $q$ whole ball rows, by contrast,
has dimension $qs+1$ and maximum exposed-face dimension $1+(q-1)s$
by Theorem~\ref{thm:whole-row-caps}. Its exact ambient parameter is
$q+1$, as the vector-row specialization of
Theorem~\ref{thm:shared-spectral-parameter}. For $k$ rows in $g$
groups the ambient parameter is $k+g$, while fixing each scale to
one gives exactly $-\sum_a\log(1-\norm{x_a}^2)$ with intrinsic
parameter $k$, independently of grouping. Thus factor sharing,
exposed-face dimension, and the barrier parameter on the actual slice
measure different properties of a formulation.

The correspondence between full slack factorizations and cone lifts is
classical~\cite{GPT2013}. Our bounds are the direct-sum
inequalities~\eqref{eq:face-packing} across contact cylinders and the
global budgets~\eqref{eq:face-global-budgets}; both assume a global
$C^1$ choice of factors for every labelled row. The barriers in the
examples are classical spectral and PSD-completion barriers. Neither the
face bounds nor a small ambient count alone implies an unrestricted
algorithmic speedup.
